\documentclass{article}
\usepackage[UTF8]{ctex}
\usepackage{amsmath,amssymb,amsfonts}
\usepackage{enumitem}
\usepackage{geometry}
\geometry{a4paper, margin=1in}
% === 错题本排版优化（不参与题目标识） ===
\setlength{\parindent}{0pt}
\setlength{\parskip}{0.35em plus 0.10em minus 0.08em}
\setlength{\emergencystretch}{3em}
\setlist{itemsep=0.25em, topsep=0.40em, parsep=0pt}
\allowdisplaybreaks[2]
% === 错题本排版优化结束 ===
\begin{document}
\section*{福州大学2024年数学分析真题}
\section*{116．福州大学2024年数学分析试卷}

\subsection{题目1}
一、（20 分）记 \(a_{n}=\int_{0}^{\frac{\pi}{4}} \tan ^{n} x \mathrm{~d} x, n \geq 1\) ．

\begin{itemize}
  \item[（1）] 证明：\(\left\{a_{n}\right\}\) 收敛。
  \item[（2）] 证明：\(n \geq 2\) 时，有 \(\frac{1}{2(n+1)} \leq a_{n} \leq \frac{1}{2(n-1)}\) ．
\end{itemize}

\subsection{题目2}
二．（20 分）设 \(f(x)\) 在 \([0,1]\) 上二阶可导，\(f^{\prime}(0)=f^{\prime}(1)=0\) ．证明：存在 \(\xi \in(0,1)\) ，满足

\[
\left|f^{\prime \prime}(\xi)\right| \geq 4|f(1)-f(0)|
\]

\subsection{题目3}
三．（20分）分析级数 \(\frac{1}{2}-\frac{1}{5}+\frac{1}{8}-\frac{1}{11}+\cdots\) 收敛吗？若收敛，求级数的和．

\subsection{题目4}
四．（20 分）设 \(f(x), g(x)\) 在 \([a, b]\) 上连续，证明：

\[
\left(\int_{a}^{b} f(x) g(x) \mathrm{d} x\right)^{2} \leq \int_{a}^{b} f^{2}(x) \mathrm{d} x \int_{a}^{b} g^{2}(x) \mathrm{d} x
\]

取等号当且仅当存在常数 \(\alpha, \beta\) 使得 \(\alpha f(x)=\beta g(x)\) ．

\subsection{题目5}
五．（20分）计算 \(I=\int_{L}-y \mathrm{~d} x+x \mathrm{~d} y, L: x^{3}+y^{3}=x y, x \geq 0, y \geq 0\) ，取逆时针方向．

\subsection{题目6}
六．（25 分）计算单位时间内某流体 \(v=y z j\)（密度为 1 ）通过曲面 \(S\) 的流量，\(S\) 为 \(x^{2}+\frac{y^{2}}{4}+z^{2}=1\) 的上半椭球面，取上侧．

\subsection{题目7}
七．（25 分）计算 \(I=\int_{0}^{+\infty} \frac{e^{-x}-e^{-2 x}}{x} \cos x \mathrm{~d} x\) ．
\end{document}
